\documentclass[11pt]{article}
% Standalone preamble for web-only supplementary notes.
\usepackage[T1]{fontenc}
\usepackage{amsmath,amssymb,amsthm}
\usepackage{hyperref}
\usepackage{cleveref}
\DeclareMathOperator{\Map}{Map}
\DeclareMathOperator{\Fun}{Fun}
\newcommand{\iso}{\xrightarrow{\sim}}
\newtheorem{lemma}{Supplementary Lemma}
\newtheorem{corollary}[lemma]{Supplementary Corollary}
\crefname{lemma}{Supplementary Lemma}{Supplementary Lemmas}
\Crefname{lemma}{Supplementary Lemma}{Supplementary Lemmas}
\crefname{corollary}{Supplementary Corollary}{Supplementary Corollaries}
\Crefname{corollary}{Supplementary Corollary}{Supplementary Corollaries}

\begin{document}
\section*{Identifications on coproducts}

The \href{initial-categories-and-coproducts.html}{coproduct universal property} is stated using
mapping animae. To glue identifications of functors, we need the
corresponding assertion for their identification animae, with the
specified restriction functor. The passage between these formulations
uses the decoding square from
\href{supplementary-decoding-and-restriction.html}{Decoding mapping animae and restriction}.
All categories and functors below are absolute. We write $\star$ for
whiskering, $D(u)$ for the functor decoded from a point $u$ of a mapping
anima, and $D_{u,v}$ for the decoding functor on identification animae.

\begin{lemma}
\label[lemma]{lem:Supplementary_Paired_Observation}
Let $F\colon X\to Y$ and $G\colon X\to Z$ be functors such that
$(F,G)\colon X\to Y\times Z$ is an equivalence. For $u,v\colon T\to X$,
the functor
\[
 (F\star-,G\star-)\colon(u=v)\longrightarrow
       (F\circ u=F\circ v)\times(G\circ u=G\circ v)
\]
is an equivalence.
\end{lemma}
\begin{proof}[Proof idea]
Put $H=(F,G)$. Postwhiskering by $H$ is an equivalence on identification
animae. Apply the product comparison to its target $(H\circ u=H\circ v)$.
The two resulting components are successive postwhiskerings by $H$ and
by a projection. Associator compatibility and the projection
identifications $\pi_1\circ H=F$ and $\pi_2\circ H=G$ compare these
components with postwhiskering by $F$ and $G$.

Each comparison is a square with possibly different endpoints. Changing
endpoints by the specified comparisons gives equivalences of the target
identification animae. Pair the two squares and transfer equivalence
across them. Thus it is the displayed paired postwhiskering functor
itself that is an equivalence.
\end{proof}

\section*{Restricting to the two summands}

\begin{lemma}
\label[lemma]{lem:Supplementary_Coproduct_Restriction}
For $h,k\colon C\amalg D\to E$, restriction along the inclusions
$i_C\colon C\to C\amalg D$ and $i_D\colon D\to C\amalg D$ is an
equivalence
\[
 R_{h,k}\colon(h=k)\longrightarrow
 (h\circ i_C=k\circ i_C)\times(h\circ i_D=k\circ i_D),
 \qquad R_{h,k}=(-\star i_C,-\star i_D).
\]
\end{lemma}
\begin{proof}[Proof idea]
First suppose $h=D(u)$ and $k=D(v)$ are decoded absolute objects of
$\Map(C\amalg D,E)$. The coproduct axiom says that
\[
 (i_C^*,i_D^*)\colon\Map(C\amalg D,E)
                  \longrightarrow\Map(C,E)\times\Map(D,E)
\]
is an equivalence. Apply \Cref{lem:Supplementary_Paired_Observation}
to this functor and to $u,v$. Decode the two target identification
animae, using that decoding acts by equivalences.

The decoding square identifies these two components with restriction
of decoded identifications, retaining the comparisons
$D(i_C^*u)=D(u)\circ i_C$ and $D(i_D^*u)=D(u)\circ i_D$, and the
corresponding comparisons for $v$. Pair the two squares, change endpoints,
and cancel the equivalence $D_{u,v}$. This proves that
$R_{D(u),D(v)}$ is an equivalence.

For general $h,k$, take their names in the mapping anima. The decoding
comparisons give $D(\operatorname{name}(h))=h$ and
$D(\operatorname{name}(k))=k$. Restriction commutes with changing these
endpoints: prewhisker the endpoint-change square by each inclusion.
Transport the decoded result across the resulting paired square to
obtain the asserted equivalence $R_{h,k}$.
\end{proof}

\begin{corollary}
\label[corollary]{cor:Supplementary_Coproduct_Lift}
Given identifications $\sigma\colon h\circ i_C=k\circ i_C$ and
$\tau\colon h\circ i_D=k\circ i_D$, there is an identification
$\lambda\colon h=k$ equipped with comparisons
\[
 \lambda\star i_C=\sigma,
 \qquad \lambda\star i_D=\tau.
\]
\end{corollary}
\begin{proof}[Proof idea]
Lift the absolute object $(\sigma,\tau)$ along the equivalence $R_{h,k}$
of \Cref{lem:Supplementary_Coproduct_Restriction}. This produces both
$\lambda$ and a comparison $R_{h,k}(\lambda)=(\sigma,\tau)$.
Postcompose with each projection and apply the pairing projection
comparisons. These are the two specified restriction comparisons.
\end{proof}

Keeping these comparisons makes the construction usable when gluing
data with prescribed restrictions. The argument uses the mapping-anima
universal property of $C\amalg D$; it does not require the coproduct
itself to be an anima.
\end{document}
